\begin{section_1}

\subsection{Свойство 1}
\textit{Утверждение:} Если $c \in M$, то и $\overline{c} \in M$, где $M$ — множество Мандельброта, а $\overline{c}$ — комплексно-сопряжённое к $c$.


\textit{Доказательство:}

\begin{enumerate}
    \item Напомним, что множество Мандельброта $M$ — это множество таких $c \in \mathbb{C}$, что последовательность
    \[
        z_0 = 0,\qquad z_{n+1} = z_n^2 + c
    \]
    остаётся ограниченной, то есть существует $R > 0$, что $\forall n$ выполнено $|z_n| < R$.
    \item Рассмотрим сопряжённое комплексное число $\overline{c}$. Построим последовательность $(w_n)$ с $w_0=0$:
    \[
        w_{n+1} = w_n^2 + \overline{c}
    \]
    Докажем, что $w_n = \overline{z_n}$ для любого $n$ по индукции:
    \begin{itemize}
        \item \textit{База:} $n=0$, $w_0 = 0 = \overline{0} = \overline{z_0}$.
        \item \textit{Переход:} Пусть $w_n = \overline{z_n}$ для некоторого $n$. Тогда:
        \begin{align*}
            w_{n+1} &= w_n^2 + \overline{c} \\
                    &= (\overline{z_n})^2 + \overline{c} \\
                    &= \overline{z_n^2} + \overline{c} \\
                    &= \overline{z_n^2 + c} \\
                    &= \overline{z_{n+1}}
        \end{align*}
        Следовательно, $w_{n+1} = \overline{z_{n+1}}$.
    \end{itemize}
    \item Таким образом, по индукции $w_n = \overline{z_n}$ для любого $n$.
    \item Значит, $|w_n| = |\overline{z_n}| = |z_n|$ (модуль числа не меняется при сопряжении!). Следовательно, ограниченность последовательности $(z_n)$ эквивалентна ограниченности $(w_n)$.
    \item Итак, $c \in M \iff \overline{c} \in M$, т.е. $M$ симметрично относительно вещественной оси.
\end{enumerate}

\vspace{1em}

\textbf{Почему последовательность $z_n$ уходит в бесконечность при $|z_n| > 2$:}

Если на каком-то шаге $m$ выполнено $|z_m| > 2$, то для любого $n \geq m$:
\begin{align*}
    |z_{n+1}| &= |z_n^2 + c| \\
              &\geq \left|\,|z_n|^2 - |c|\,\right|.
\end{align*}
Поскольку $|z_n| > 2$, то $|z_n|^2 > 4$. Если выбрать $|z_n|$ настолько большим, что $|z_n|^2 > 2|c|$, то
\[
    |z_{n+1}| \geq |z_n|^2 - |c| > \frac{1}{2}|z_n|^2,
\]
то есть $|z_{n+1}|$ строго больше, чем $|z_n|$, и перестаёт быть ограниченным сверху. Следовательно, начиная с этого момента, последовательность модулей $|z_n|$ при каждом переходе хотя бы удваивается, стремясь к бесконечности. Поэтому если $|z_n|$ хоть раз превысит 2, уже далее орбита обязательно покинет любой круг конечного радиуса и не вернётся внутрь него.


\vspace{1.5em}
\subsection{Свойство 2}


\textit{Утверждение:} Если $|c| > 2$, то $c$ не принадлежит множеству Мандельброта.

\textit{Доказательство:}

\begin{enumerate}
    \item Пусть $|c| > 2$. Построим последовательность $(z_n)$:
    \[
        z_0 = 0,\qquad z_{n+1} = z_n^2 + c
    \]
    \item Уже на первом шаге $z_1 = c$, поэтому $|z_1| = |c| > 2$.
    \item Покажем, что если на каком-либо шаге $|z_n| > 2$, то на следующем шаге обязательно $|z_{n+1}| > |z_n|$.
    \item Заметим, что для любого $a, b \in \mathbb{C}$:
    \[
        |a + b| \geq |\,|a| - |b|\,|
    \]
    Тогда, учитывая $z_{n+1} = z_n^2 + c$:
 \begin{align*}
    |z_{n+1}| &= |z_n^2 + c| \\
              &\geq \left||z_n|^2 - |c|\right|
\end{align*}
    \item Для $|z_n| > 2$, рассмотрим $|z_n|^2 - |c|$: ведь $|z_n| > |c| > 2$, т.е. $|z_n|^2 - |c| > |c|^2 - |c| > 2^2 - |c| \geq 4 - |c|$. Но $|c| > 2$, поэтому $|z_{n+1}| > 2$ и, более того, $|z_{n+1}| > |z_n|$.
    \item Следовательно, начиная с $z_1$, все последующие члены последовательности будут по модулю строго больше $2$ и расти неограниченно, то есть последовательность уходит на бесконечность.
    \item Значит, при $|c| > 2$ последовательность не может остаться ограниченной, то есть $c \notin M$.
\end{enumerate}

\end{section_1}